% AI 提示：附录包含支撑材料文件列表 + 全部源代码
%   - 表格：\csvtable{data/xxx.csv}{标题}{tab:xxx}（CSV 必须 UTF-8）
%   - 代码：\lstinputlisting[style=Python]{code/q1.py}（已放在 code/）

\clearpage
\phantomsection
\addcontentsline{toc}{section}{附录}
\section*{附录}

\subsection*{支撑材料文件列表}
本文提交的支撑材料文件列表如下：

\begin{table}[H]
  \centering
  \caption{支撑材料文件列表}
  \label{tab:filelist}
  \begin{tabularx}{\textwidth}{@{}>{\centering}p{4cm}>{\centering\arraybackslash}X@{}}
    \toprule
    文件名 & 功能描述 \\
    \midrule
    q1.py & 问题一的求解程序 \\
    q2.py & 问题二的求解程序 \\
    q3.py & 问题三/四的求解程序 \\
    data.xlsx & 本文使用的原始数据与处理结果 \\
    \bottomrule
  \end{tabularx}
\end{table}

\subsection*{中间结果表（CSV 自动排版示例）}
\csvtable{data/示例结果表.csv}{各模型性能对比（示例）}{tab:csvdemo}

\subsection*{核心源代码}

\subsubsection*{问题一求解程序（q1.py）}
\lstinputlisting[style=Python]{code/q1.py}

\subsubsection*{问题二求解程序（q2.py）}
\lstinputlisting[style=Python]{code/q2.py}

\subsubsection*{问题三/四求解程序（q3.py）}
\lstinputlisting[style=Python]{code/q3.py}
